\scanpage{14}
ne dépend que du Module stratifié $\underline V=V^{q_0}(M^\bullet)$ et de \deleted{\unsure{$q_0$}} $q_0$, mais pas du complexe différentiel particulier choisi pour ``résoudre'' $\underline V$.

6. Or, \deleted{\unsure{si $X$ est lisse sur $S$} \uline{\unsure{et de}} \uline{\unsure{car. $0$}} \unsure{(absolue)}} si on part d'un Module stratifié $\underline V$ sur $X$ rel / $S$, on peut tjrs lui associer de façon naturelle un \uline{complexe de De Rham}, dont l'objet gradué sous-jacent est
\[
\underline\Omega^\bullet_{X/S}\otimes\underline V=\Omega^\bullet(\underline V),
\]
et on vérifie, \uline{lorsque} \inserted{$X$ est lisse sur $S$ et,} \uline{$S$ est de car. nulle}, que ce complexe est bien une ``résolution différentielle'' de $\underline V$.
\margin{\figure{\begin{tikzpicture}
\node[draw,ellipse,rotate=25,align=center,inner sep=3pt] {``Lemme de\\Poincaré formel''};
\draw (-.8,.4)--(-.8,1.1);\draw (-.73,.4)--(-.73,1.1);\draw (-.66,.4)--(-.66,1.1);
\end{tikzpicture}}}
Dans ce cas, on trouve donc une interprétation \unsure{univoque} de la
\margin{\figure{\begin{tikzpicture}
\node[draw,ellipse,inner sep=4pt] {\unsure{différentielles|différentielle}};
\end{tikzpicture}}}
\margin{13}
